\documentclass[11pt]{article}
\usepackage[margin=2.5cm]{geometry}
\usepackage{booktabs}

\title{A Short Note on Stars}
\author{A. Student}
\date{Draft 2}

\begin{document}
\maketitle

\begin{abstract}
Stars are giant balls of hot gas held together by their own gravity. In this note
we describe how stars produce energy by nuclear fusion, how their lifetime depends
on their mass, and how their final state (white dwarf, neutron star or black hole)
is set by the mass they have left at the end of their lives.
\end{abstract}

\section{Introduction}

The Sun is an ordinary, medium-sized star and the closest star to the Earth.
Light from the Sun takes about 8 minutes to reach us. Stars produce their energy by
nuclear fusion in their cores, where hydrogen nuclei fuse into helium and release
energy~\cite{carroll}.

\section{Stellar evolution}

A star spends most of its life on the main sequence, where the pressure of the hot
gas balances gravity. Massive stars live for a much shorter time than small stars,
because they use up their fuel faster. Table~\ref{tab:stars} lists three typical
stars.

\begin{table}[h]
\centering
\caption{Typical main-sequence stars. Lifetimes are order-of-magnitude estimates.}
\label{tab:stars}
\begin{tabular}{lccc}
\toprule
Star type & Mass (Sun $=1$) & Surface temperature (K) & Lifetime (years) \\
\midrule
Red dwarf   & 0.2 & 3\,000  & $10^{12}$ \\
Sun-like    & 1   & 5\,800  & $10^{10}$ \\
Blue giant  & 20  & 30\,000 & $10^{7}$ \\
\bottomrule
\end{tabular}
\end{table}

When a star runs out of hydrogen in its core it leaves the main sequence. Its final
state depends on its mass: stars like the Sun end as white dwarfs, while the most
massive stars collapse into neutron stars or black holes.

\begin{thebibliography}{1}
\bibitem{carroll} B.~W. Carroll and D.~A. Ostlie, \emph{An Introduction to Modern
Astrophysics}, 2nd ed., Cambridge University Press, 2017.
\end{thebibliography}

\end{document}
